\chapter{Likelihood, Estimation and Bayes}
\label{ch:likelihood-bayes}

Chapters 12 and 13 fitted lines and logistic curves by minimising a loss: squared
error for regression, cross-entropy for classification. Those losses were handed to
you. This chapter explains where they come from. The answer is one idea,
\emph{likelihood}, and it also tells us how precise a fitted parameter is and how to
fold prior knowledge into a fit. By the end, ridge and lasso regression will turn out to
be the same idea again, with a prior attached.

\section{Estimators and how they miss}
\label{sec:lik-estimators}

\subsection{An estimator is a recipe}

Five API calls to a new service took these times in milliseconds:
\[
  4,\quad 7,\quad 5,\quad 9,\quad 5 .
\]
We want the true mean latency $\mu$ and the true variance $\sigma^2$ of this service.
Two natural guesses are the sample mean and the sample variance:
\[
  \xbar = \frac{4+7+5+9+5}{5} = 6, \qquad
  \sum (x_i-\xbar)^2 = 4+1+1+9+1 = 16 .
\]
Dividing $16$ by $n-1=4$ gives $s^2=4$ (Chapter 2). Dividing by $n=5$ gives $3.2$.
Both are reasonable recipes. Which is better? To answer that we need to say what
``better'' means for a recipe, not for one number.

\begin{definition}
An \textbf{estimator} $\hat\theta$ is a function of the data that is used to guess a
parameter $\theta$. Picture a machine: data go in, one number comes out. The number it
produces on the data you actually have is the \textbf{estimate}.
\end{definition}

Because the data are random, the estimator is a random variable. If we reran the
experiment, we would get five different latencies and a different $\xbar$. The spread
of these hypothetical outputs is the \textbf{sampling distribution} of the estimator
(Chapter 6): the histogram you would see if you repeated the whole study many times.

\begin{figure}[H]
\centering
\begin{tikzpicture}[
  box/.style={draw, rounded corners, minimum width=2.1cm, minimum height=0.8cm, align=center, font=\small},
  >=Stealth]
\node[box, fill=histfill!30] (pop) at (0,0) {population\\ parameter $\theta$};
\node[box] (s1) at (3.4,1.3) {sample 1};
\node[box] (s2) at (3.4,0) {sample 2};
\node[box] (s3) at (3.4,-1.3) {sample 3};
\node[box, fill=quotientorange!20] (e1) at (6.6,1.3) {$\hat\theta_1 = 6.0$};
\node[box, fill=quotientorange!20] (e2) at (6.6,0) {$\hat\theta_2 = 5.4$};
\node[box, fill=quotientorange!20] (e3) at (6.6,-1.3) {$\hat\theta_3 = 6.8$};
\draw[->] (pop) -- (s1); \draw[->] (pop) -- (s2); \draw[->] (pop) -- (s3);
\draw[->] (s1) -- node[above, font=\scriptsize] {recipe} (e1);
\draw[->] (s2) -- node[above, font=\scriptsize] {recipe} (e2);
\draw[->] (s3) -- node[above, font=\scriptsize] {recipe} (e3);
\begin{scope}[shift={(8.6,-1.3)}]
\draw[->] (0,0) -- (3.2,0);
\fill[histfill] (0.4,0) rectangle (0.8,0.5);
\fill[histfill] (0.8,0) rectangle (1.2,1.3);
\fill[histfill] (1.2,0) rectangle (1.6,2.1);
\fill[histfill] (1.6,0) rectangle (2.0,1.6);
\fill[histfill] (2.0,0) rectangle (2.4,0.7);
\fill[histfill] (2.4,0) rectangle (2.8,0.2);
\node[font=\scriptsize, align=center] at (1.6,2.55) {sampling\\ distribution of $\hat\theta$};
\end{scope}
\draw[->, dashed] (e2.east) -- (8.85,0);
\end{tikzpicture}
\caption{The same recipe applied to different samples gives different estimates. Judging
an estimator means judging the whole histogram of its outputs, not one number.}
\label{fig:lik-estimator-machine}
\end{figure}

\subsection{Bias, variance and mean squared error}

Two things can go wrong with a recipe. Its outputs can be centred in the wrong place,
or they can be scattered widely.

\begin{definition}
The \textbf{bias} of $\hat\theta$ is $\Bias(\hat\theta) = \E[\hat\theta] - \theta$: how
far the centre of the sampling distribution sits from the truth. The \textbf{variance}
$\Var(\hat\theta)$ measures how scattered the outputs are around their own centre. The
\textbf{mean squared error} is $\MSE(\hat\theta) = \E[(\hat\theta-\theta)^2]$: the
average squared distance from the truth.
\end{definition}

A dartboard makes the two failure modes visible. The bullseye is $\theta$; each dart is
the estimate from one repeated study.

\begin{figure}[H]
\centering
\begin{tikzpicture}[scale=0.82]
\foreach \cx/\cy/\lab in {0/0/{low bias, low variance}, 5.2/0/{low bias, high variance},
                         0/-4.4/{high bias, low variance}, 5.2/-4.4/{high bias, high variance}} {
  \begin{scope}[shift={(\cx,\cy)}]
  \fill[histfill!25] (0,0) circle (1.6);
  \fill[histfill!50] (0,0) circle (1.07);
  \fill[histfill!80] (0,0) circle (0.53);
  \draw[gray] (0,0) circle (1.6) (0,0) circle (1.07) (0,0) circle (0.53);
  \fill[black] (0,0) circle (0.05);
  \node[font=\small, anchor=north] at (0,-1.7) {\lab};
  \end{scope}
}
\foreach \x/\y in {0.15/0.1, -0.1/0.2, 0.2/-0.15, -0.2/-0.05, 0.05/-0.2, -0.05/0.05, 0.25/0.05, -0.15/-0.2}
  \fill[densityred] (\x,\y) circle (0.09);
\foreach \x/\y in {0.7/0.5, -0.6/0.8, 0.9/-0.4, -0.8/-0.3, 0.1/-0.9, -0.3/0.2, 0.4/1.0, -0.4/-0.9}
  \fill[densityred] (5.2+\x,\y) circle (0.09);
\foreach \x/\y in {0.15/0.1, -0.1/0.2, 0.2/-0.15, -0.2/-0.05, 0.05/-0.2, -0.05/0.05, 0.25/0.05, -0.15/-0.2}
  \fill[densityred] (0.8+\x,-4.4+0.6+\y) circle (0.09);
\foreach \x/\y in {0.7/0.5, -0.6/0.8, 0.9/-0.4, -0.8/-0.3, 0.1/-0.9, -0.3/0.2, 0.4/1.0, -0.4/-0.9}
  \fill[densityred] (5.2+0.8+0.7*\x,-4.4+0.6+0.7*\y) circle (0.09);
\end{tikzpicture}
\caption{Bias is where the cloud of darts is centred; variance is how wide the cloud is.
The top-left estimator is the goal. The bottom-left one is precise but consistently
wrong, and no amount of repetition reveals that from the darts alone.}
\label{fig:lik-dartboard}
\end{figure}

The mean squared error charges for both failures, and it splits cleanly into the two.

\begin{theorem}[Bias--variance decomposition]
For any estimator, $\MSE(\hat\theta) = \Bias(\hat\theta)^2 + \Var(\hat\theta)$.
\end{theorem}

\begin{whybox}
Write $m = \E[\hat\theta]$ for the centre of the darts. Add and subtract it:
\[
  \hat\theta - \theta = \underbrace{(\hat\theta - m)}_{\text{scatter}}
  + \underbrace{(m - \theta)}_{\text{bias}} .
\]
Square and take expectations:
\[
  \E[(\hat\theta-\theta)^2] = \E[(\hat\theta-m)^2] + 2(m-\theta)\,\E[\hat\theta-m]
  + (m-\theta)^2 .
\]
The middle term vanishes because $\E[\hat\theta - m] = m - m = 0$: scatter around the
centre averages to zero. The first term is $\Var(\hat\theta)$ and the last is
$\Bias^2$. It is Pythagoras: the error vector has a ``random'' leg and a ``systematic''
leg at right angles.
\end{whybox}

\begin{recoverybox}
Pause and think: an estimator always outputs the number $5$, whatever the data. What
are its variance and its MSE? (Variance $0$; MSE $=(5-\theta)^2$. Superb if
$\theta=5$, useless otherwise. Low variance alone is not a virtue.)
\end{recoverybox}

\subsection{Unbiased is not the same as best}

Return to the variance of the latencies. For Normal data with $n$ points, the sum of
squared deviations $SS=\sum(x_i-\xbar)^2$ satisfies $\E[SS]=(n-1)\sigma^2$ and
$\Var(SS) = 2(n-1)\sigma^4$ (see Probability From the Ground Up, Ch.~30). So dividing
$SS$ by some number $d$ gives
\[
  \Bias = \frac{n-1}{d}\sigma^2 - \sigma^2, \qquad
  \Var = \frac{2(n-1)}{d^2}\sigma^4 .
\]
With $n=5$ and $\sigma^2=1$:

\begin{center}
\begin{tabular}{lccc}
\toprule
divisor $d$ & bias$^2$ & variance & MSE \\
\midrule
$n-1=4$ (Bessel, Chapter 2) & $0$ & $0.500$ & $0.500$ \\
$n=5$ & $0.040$ & $0.320$ & $0.360$ \\
$n+1=6$ & $0.111$ & $0.222$ & $0.333$ \\
\bottomrule
\end{tabular}
\end{center}

\begin{figure}[H]
\centering
\begin{tikzpicture}
\begin{axis}[
  ybar stacked, width=9cm, height=5.6cm, bar width=1.1cm,
  symbolic x coords={d=4,d=5,d=6}, xtick=data,
  ymin=0, ymax=0.62, ylabel={error (units of $\sigma^4$)},
  xlabel={divisor of $SS$ with $n=5$},
  legend style={at={(0.98,0.97)}, anchor=north east, font=\small},
  axis lines=left, enlarge x limits=0.3,
]
\addplot[fill=quotientorange!70, draw=quotientorange] coordinates {(d=4,0) (d=5,0.04) (d=6,0.111)};
\addlegendentry{bias$^2$}
\addplot[fill=histfill, draw=distblue] coordinates {(d=4,0.5) (d=5,0.32) (d=6,0.222)};
\addlegendentry{variance}
\end{axis}
\end{tikzpicture}
\caption{Dividing by a larger number shrinks the estimate. That adds a little bias but
removes more variance, so total MSE falls. The unbiased choice $d=n-1$ has the
largest MSE of the three.}
\label{fig:lik-divisor-mse}
\end{figure}

Unbiasedness is a property of the centre of the darts. It says nothing about how close
any single dart lands. Machine learning leans on exactly this trade: regularisation
(Section~\ref{sec:lik-map}) deliberately adds bias to cut variance.

\subsection{Consistency}

An estimator is \textbf{consistent} if it closes in on the truth as the sample grows:
for any tolerance $\eps>0$, $\P(\abs{\hat\theta_n-\theta}>\eps)\to 0$ as $n\to\infty$.
In pictures, the sampling distribution piles up into a spike at $\theta$. A simple
sufficient condition is that both bias and variance go to zero, because then MSE goes
to zero.

\begin{figure}[H]
\centering
\begin{tikzpicture}
\begin{axis}[
  width=11cm, height=5.6cm, domain=0.2:1, samples=200,
  xlabel={$\hat p$, the proportion of heads}, ylabel={density},
  ymin=0, ymax=13.5, axis lines=left,
  legend style={at={(0.02,0.97)}, anchor=north west, font=\small},
]
\addplot[densityred, thick] {exp(-(x-0.7)^2/(2*0.021))/sqrt(2*pi*0.021)};
\addlegendentry{$n=10$}
\addplot[quotientorange, thick] {exp(-(x-0.7)^2/(2*0.0042))/sqrt(2*pi*0.0042)};
\addlegendentry{$n=50$}
\addplot[distblue, thick] {exp(-(x-0.7)^2/(2*0.00105))/sqrt(2*pi*0.00105)};
\addlegendentry{$n=200$}
\draw[gray, dashed] (axis cs:0.7,0) -- (axis cs:0.7,13.2);
\node[gray, anchor=west, font=\small] at (axis cs:0.72,12.6) {$p=0.7$};
\end{axis}
\end{tikzpicture}
\caption{Sampling distribution of the proportion $\hat p$ when the true $p=0.7$
(Normal approximation, SD $\sqrt{0.21/n}$). The curve tightens around the truth as
$n$ grows: $\hat p$ is consistent.}
\label{fig:lik-consistency}
\end{figure}

Both variance recipes above are consistent: their bias ($0$ or $-\sigma^2/n$) and their
variance both vanish as $n$ grows. At large $n$ the choice of divisor stops mattering.

\subsection{Method of moments, quickly}

The oldest recipe for building estimators: match sample averages to their theoretical
formulas and solve. A \textbf{moment} is an average of a power, such as $\E[X]$ or
$\E[X^2]$.

\begin{example}
A job's runtime in seconds is uniform between $0$ and an unknown maximum $\theta$.
Five runs took $3.1$, $0.4$, $5.2$, $1.9$ and $4.4$, so $\xbar = 3.0$. Theory says
$\E[X]=\theta/2$. Setting $\theta/2 = 3.0$ gives $\hat\theta_{\text{MoM}} = 6.0$.
\end{example}

Method of moments is quick and usually consistent, but it can be wasteful or even
silly: if one run had taken $6.5$ seconds while $\xbar$ stayed at $3.0$, it would
still say $\theta = 6.0$, a value the data have already ruled out. The method we turn to
next never makes that mistake.

\section{Likelihood}
\label{sec:lik-likelihood}

\subsection{Turning the probability around}

A coin is flipped $10$ times and lands heads $7$ times. Which values of $p$, the chance
of heads, are plausible?

For any candidate $p$ we can compute the probability of what we saw. For one specific
sequence with $7$ heads and $3$ tails, that probability is $p^7(1-p)^3$:
\begin{align*}
  p=0.5:&\quad 0.5^7\,0.5^3 = 0.000977,\\
  p=0.7:&\quad 0.7^7\,0.3^3 = 0.002224 .
\end{align*}
The data are about $2.28$ times more probable under $p=0.7$ than under $p=0.5$. Do this
for every $p$ and plot.

\begin{definition}
The \textbf{likelihood function} is $L(\theta) = p(\text{data}\given\theta)$, read as a
function of $\theta$ with the data held fixed. It scores each candidate parameter by
how well it predicts the data you already have.
\end{definition}

\begin{figure}[H]
\centering
\begin{tikzpicture}
\begin{axis}[
  width=11cm, height=5.8cm, domain=0:1, samples=200,
  xlabel={$p$}, ylabel={$L(p)=p^7(1-p)^3$},
  ymin=0, ymax=0.0026, axis lines=left,
  scaled y ticks=false, yticklabel style={/pgf/number format/fixed, /pgf/number format/precision=4},
]
\addplot[name path=L, distblue, very thick] {x^7*(1-x)^3};
\path[name path=axis] (axis cs:0,0) -- (axis cs:1,0);
\addplot[histfill!40] fill between[of=L and axis];
\draw[densityred, dashed] (axis cs:0.7,0) -- (axis cs:0.7,0.002224);
\fill[densityred] (axis cs:0.7,0.002224) circle (2pt);
\node[densityred, anchor=south west, font=\small] at (axis cs:0.71,0.00223) {peak at $\hat p = 0.7$};
\draw[gray, dashed] (axis cs:0.5,0) -- (axis cs:0.5,0.000977);
\fill[gray] (axis cs:0.5,0.000977) circle (2pt);
\node[gray, anchor=east, font=\small] at (axis cs:0.49,0.00105) {$L(0.5)=0.00098$};
\end{axis}
\end{tikzpicture}
\caption{Likelihood of $p$ after $7$ heads in $10$ flips. It peaks at $0.7$, is
still sizeable at $0.5$, and is essentially zero below $0.2$. The width of the hump
shows how uncertain $p$ still is.}
\label{fig:lik-coin-likelihood}
\end{figure}

\begin{trapbox}
The likelihood is not a probability distribution over $\theta$. It does not integrate
to $1$ (the area under Figure~\ref{fig:lik-coin-likelihood} is about $0.00076$), and
``$L(0.7)=0.0022$'' does not mean ``$p=0.7$ has probability $0.0022$''. Only ratios
of likelihoods carry meaning. Turning likelihood into a distribution over $\theta$
requires a prior, which is the Bayesian step in Section~\ref{sec:lik-bayes}.
\end{trapbox}

The \textbf{maximum likelihood estimator} (MLE) is the parameter that makes the data
most probable:
\[
  \hat\theta_{\text{MLE}} = \argmax_\theta L(\theta).
\]

\subsection{Logs make it easy}

For independent observations the likelihood is a product,
$L(\theta)=\prod_i p(x_i\given\theta)$. Products of many small numbers underflow on a
computer and are awkward to differentiate. Take logs: the \textbf{log-likelihood}
$\ell(\theta) = \log L(\theta) = \sum_i \log p(x_i\given\theta)$ turns the product into a
sum. The logarithm is increasing, so it peaks at the same $\theta$.

\begin{figure}[H]
\centering
\begin{tikzpicture}
\begin{axis}[
  width=11cm, height=5.8cm, domain=0.05:0.97, samples=200,
  xlabel={$p$}, ylabel={$\ell(p) = 7\log p + 3\log(1-p)$},
  ymin=-22, ymax=-3, axis lines=left,
]
\addplot[distblue, very thick] {7*ln(x)+3*ln(1-x)};
\draw[densityred, thick] (axis cs:0.5,-6.1086) -- (axis cs:0.9,-6.1086);
\fill[densityred] (axis cs:0.7,-6.1086) circle (2pt);
\node[densityred, anchor=south, font=\small] at (axis cs:0.7,-5.6) {slope $0$ at $\hat p=0.7$, $\ell=-6.11$};
\fill[gray] (axis cs:0.5,-6.9315) circle (2pt);
\node[gray, anchor=north west, font=\small] at (axis cs:0.51,-7.4) {$\ell(0.5)=-6.93$};
\end{axis}
\end{tikzpicture}
\caption{The log-likelihood has the same peak as the likelihood but is a smooth hill
that is easy to differentiate. The MLE is where the tangent is flat.}
\label{fig:lik-coin-loglik}
\end{figure}

\subsection{MLE for a Bernoulli proportion}

With $k$ successes in $n$ trials, $\ell(p) = k\log p + (n-k)\log(1-p)$. Set the
derivative to zero:
\[
  \ell'(p) = \frac{k}{p} - \frac{n-k}{1-p} = 0
  \;\Longrightarrow\; k(1-p) = (n-k)p
  \;\Longrightarrow\; \hat p = \frac{k}{n}.
\]
The MLE is the observed proportion, $7/10$. Here the method only confirms common sense;
the next two examples show it producing answers that are less obvious.

\subsection{MLE for a Normal: why it divides by $n$}

Take $x_1,\dots,x_n \iid \Normal(\mu,\sigma^2)$. Each point contributes
$\log p(x_i) = -\tfrac12\log(2\pi\sigma^2) - (x_i-\mu)^2/(2\sigma^2)$, so
\[
  \ell(\mu,\sigma^2) = -\frac n2\log(2\pi\sigma^2)
  - \frac{1}{2\sigma^2}\sum_{i=1}^n (x_i-\mu)^2 .
\]
Maximising over $\mu$ only involves the sum of squares, which is smallest at
$\hat\mu = \xbar$. Plug that in and differentiate in $v=\sigma^2$:
\[
  \frac{\partial\ell}{\partial v} = -\frac{n}{2v} + \frac{SS}{2v^2} = 0
  \;\Longrightarrow\;
  \hat\sigma^2_{\text{MLE}} = \frac{1}{n}\sum_{i=1}^n (x_i-\xbar)^2 .
\]
For the latencies, $\hat\mu = 6$ and $\hat\sigma^2 = 16/5 = 3.2$, not $s^2 = 4$.

\begin{whybox}
The MLE divides by $n$ because it plugs in $\xbar$ as if it were the true $\mu$. But
$\xbar$ was chosen to be as close as possible to these particular points, so deviations
from $\xbar$ are systematically smaller than deviations from $\mu$. Chapter 2 fixed this
by dividing by $n-1$. The MLE does not correct, so it is biased low by the factor
$(n-1)/n$. As Figure~\ref{fig:lik-divisor-mse} showed, that bias is not a disaster: the
MLE has lower MSE than $s^2$. For large $n$ the two agree.
\end{whybox}

\subsection{MLE for a Poisson rate}

A service logged these error counts in six consecutive hours: $2, 0, 3, 1, 4, 2$. A
Poisson model with rate $\lambda$ (the standard model for counts of rare, independent
events) gives $p(k\given\lambda) = e^{-\lambda}\lambda^k/k!$. Then
\[
  \ell(\lambda) = \sum_i\big(k_i\log\lambda - \lambda - \log k_i!\big)
  = 12\log\lambda - 6\lambda - \text{const},
\]
and $\ell'(\lambda) = 12/\lambda - 6 = 0$ gives $\hat\lambda = 2 = \xbar$.

\begin{figure}[H]
\centering
\begin{tikzpicture}
\begin{axis}[
  width=10.5cm, height=5.4cm, domain=0.6:4.5, samples=150,
  xlabel={rate $\lambda$ (errors per hour)}, ylabel={$\ell(\lambda)$},
  ymin=-17, ymax=-9, axis lines=left,
]
\addplot[distblue, very thick] {12*ln(x)-6*x-6.3561};
\fill[densityred] (axis cs:2,-10.0383) circle (2pt);
\node[densityred, anchor=south, font=\small] at (axis cs:2,-9.8) {$\hat\lambda = 2$};
\fill[gray] (axis cs:1,-12.3561) circle (2pt);
\node[gray, anchor=north west, font=\small] at (axis cs:1.05,-12.5) {$\ell(1)=-12.36$};
\fill[gray] (axis cs:3,-11.1728) circle (2pt);
\node[gray, anchor=north east, font=\small] at (axis cs:2.95,-11.4) {$\ell(3)=-11.17$};
\end{axis}
\end{tikzpicture}
\caption{Poisson log-likelihood for the counts $2,0,3,1,4,2$. It peaks at the mean
count. The hill is asymmetric: rates below the peak lose plausibility faster than rates
above it.}
\label{fig:lik-poisson}
\end{figure}

For the Uniform runtime example, the likelihood is $\theta^{-n}$ when $\theta$ is at
least the largest observation and $0$ otherwise, so the MLE is the maximum, $5.2$. It
can never propose a $\theta$ the data have ruled out. (It is biased low; multiplying by
$(n+1)/n$ fixes that, giving $6.24$.)

\begin{lstlisting}[language=Python]
import numpy as np
from scipy import stats, optimize
x = np.array([4, 7, 5, 9, 5])
print(stats.norm.fit(x))        # (6.0, 1.789): sigma_hat^2 = 3.2, divides by n
k = np.array([2, 0, 3, 1, 4, 2])
nll = lambda lam: -stats.poisson.logpmf(k, lam).sum()
print(optimize.minimize_scalar(nll, bounds=(0.01, 10), method="bounded").x)  # 2.0
\end{lstlisting}

\section{Every loss function is a likelihood}
\label{sec:lik-losses}

\subsection{Least squares is Gaussian MLE}

Regression assumes $y_i = f(x_i) + \eps_i$ with noise $\eps_i\iid\Normal(0,\sigma^2)$.
The likelihood of the observed $y$'s is the Normal density evaluated at each residual,
so the negative log-likelihood is
\[
  -\ell = \frac n2\log(2\pi\sigma^2) + \frac{1}{2\sigma^2}\sum_i \big(y_i - f(x_i)\big)^2 .
\]
For any fixed $\sigma$, the only part that depends on $f$ is the sum of squared
residuals. Maximising the likelihood is \emph{exactly} minimising squared error. The
least-squares line of Chapter 12 is the MLE under Gaussian noise, and
$\hat\sigma^2 = \text{SSE}/n$ falls out as a bonus.

\subsection{Cross-entropy is Bernoulli MLE}

A classifier outputs $\hat p_i = \P(y_i=1\given x_i)$. Each label is a coin flip with
that probability, so $p(y_i) = \hat p_i^{\,y_i}(1-\hat p_i)^{1-y_i}$ and
\[
  -\log p(y_i) = -\big[y_i\log\hat p_i + (1-y_i)\log(1-\hat p_i)\big].
\]
That is the binary cross-entropy loss from Chapter 13. A confident correct prediction
($\hat p=0.9$, $y=1$) costs $-\log 0.9 = 0.105$. A confident wrong one ($\hat p = 0.2$,
$y=1$) costs $-\log 0.2 = 1.609$. For $K$ classes the same reasoning with a categorical
distribution gives $-\log \hat p_{\text{true class}}$: softmax cross-entropy.

\subsection{Absolute error is Laplace MLE}

The \textbf{Laplace distribution} has density $\frac{1}{2b}e^{-\abs{x-\mu}/b}$: a
two-sided exponential, a sharp peak with heavier tails than the Normal. Its negative log
density is $\abs{x-\mu}/b + \log 2b$, so maximum likelihood minimises the sum of absolute
errors.

\begin{figure}[H]
\centering
\begin{tikzpicture}
\begin{axis}[
  name=dens, width=6.4cm, height=5cm, domain=-4:4, samples=161,
  title={noise density}, title style={font=\small},
  xlabel={residual}, ymin=0, ymax=0.56, axis lines=left,
  legend style={at={(0.02,0.98)}, anchor=north west, font=\scriptsize},
]
\addplot[distblue, thick] {exp(-x^2/2)/sqrt(2*pi)};
\addlegendentry{Gaussian}
\addplot[densityred, thick] {0.5*exp(-abs(x))};
\addlegendentry{Laplace}
\end{axis}
\begin{axis}[
  at={(dens.east)}, anchor=west, xshift=1.4cm,
  width=6.4cm, height=5cm, domain=-4:4, samples=161,
  title={loss $=-\log$ density (shifted)}, title style={font=\small},
  xlabel={residual}, ymin=0, ymax=8.5, axis lines=left,
  legend style={at={(0.5,0.98)}, anchor=north, font=\scriptsize},
]
\addplot[distblue, thick] {x^2/2};
\addlegendentry{squared}
\addplot[densityred, thick] {abs(x)};
\addlegendentry{absolute}
\end{axis}
\end{tikzpicture}
\caption{Each noise model on the left produces the loss on the right. Gaussian tails
die fast, so a large residual is ``impossible'' and the squared loss punishes it
heavily. Laplace tails are fatter, so big residuals cost only linearly.}
\label{fig:lik-noise-loss}
\end{figure}

Minimising $\sum(x_i-c)^2$ over a constant $c$ gives the mean; minimising
$\sum\abs{x_i-c}$ gives the median (Chapter 2). Five response times
$1,2,3,4,20$ ms make the difference concrete: the mean is $6$, the median is $3$.

\begin{figure}[H]
\centering
\begin{tikzpicture}
\draw[->] (0,0) -- (11.5,0) node[right, font=\small] {ms};
\foreach \t in {0,5,10,15,20} \draw (\t*0.55,0.08) -- (\t*0.55,-0.08) node[below, font=\scriptsize] {\t};
\foreach \v in {1,2,3,4,20} \fill[distblue] (\v*0.55,0.25) circle (0.11);
\draw[densityred, thick] (6*0.55,-0.1) -- (6*0.55,1.0);
\node[densityred, anchor=south, font=\small] at (6*0.55,1.0) {mean $6$ (Gaussian MLE)};
\draw[targetgreen, thick] (3*0.55,-0.1) -- (3*0.55,0.6);
\node[targetgreen, anchor=south east, font=\small] at (3*0.55+0.1,0.6) {median $3$ (Laplace MLE)};
\end{tikzpicture}
\caption{One slow request drags the Gaussian MLE (mean) to $6$, above four of the five
points. The Laplace MLE (median) stays with the bulk. Picking a loss is picking a belief
about how often big errors happen.}
\label{fig:lik-mean-median}
\end{figure}

\subsection{The dictionary}

\begin{center}
\begin{tabular}{llll}
\toprule
loss & noise model & best constant & used in \\
\midrule
squared error & Gaussian & mean & linear regression \\
absolute error & Laplace & median & robust regression \\
pinball loss at level $\tau$ & asymmetric Laplace & $\tau$-quantile & p99 forecasting \\
binary cross-entropy & Bernoulli & proportion & logistic regression \\
categorical cross-entropy & categorical & class frequencies & softmax classifiers \\
$\lambda - y\log\lambda$ & Poisson & mean count & count regression \\
\bottomrule
\end{tabular}
\end{center}

The \textbf{pinball loss} (Chapter 3) is a lopsided absolute error. For a residual
$r = y - c$ it charges $\tau r$ when $r \ge 0$ (you predicted too low) and $(1-\tau)\abs{r}$
when $r<0$ (too high). At $\tau = 0.99$, predicting too low costs $99$ times more per
millisecond than predicting too high, so the best constant $c$ is pushed up until only
$1\%$ of the data lie above it: the 99th percentile. At $\tau = 0.5$ it is half the
absolute error, and the best constant is the median.

\begin{connectionbox}
In ML papers the training objective is often written
$\min_\theta \E_{(x,y)\sim\text{data}}[-\log p_\theta(y\given x)]$. Read it as ``average
negative log-likelihood over the training set''. Every row of the table is this formula
with a different noise model plugged in. When a paper says ``we model the output as a
Gaussian with learned variance'', it means the loss has an extra $\log\sigma$ term and
the squared error is divided by $\sigma^2$. Chapter 15 returns to this.
\end{connectionbox}

\begin{recoverybox}
Pause and think: you want a model that predicts p99 latency, not mean latency. Should you
train it with squared error? (No. Squared error targets the mean. Use the pinball loss
at $\tau = 0.99$, whose best constant is the 99th percentile.)
\end{recoverybox}

\section{How sure is the MLE?}
\label{sec:lik-fisher}

\subsection{Curvature is information}

Look again at the log-likelihood of Figure~\ref{fig:lik-coin-loglik}. If the hill is
sharp, moving $p$ a little away from the peak costs a lot of log-likelihood, so only a
narrow range of $p$ is plausible. If the hill is flat, many values fit almost equally
well. Sharpness at the top is measured by the second derivative.

\begin{figure}[H]
\centering
\begin{tikzpicture}
\begin{axis}[
  width=11cm, height=6cm, domain=0.35:0.95, samples=200,
  xlabel={$p$}, ylabel={$\ell(p) - \ell(\hat p)$},
  ymin=-6, ymax=0.6, axis lines=left,
  legend style={at={(0.5,-0.3)}, anchor=north, font=\small}, legend columns=2,
]
\addplot[densityred, very thick] {7*ln(x/0.7)+3*ln((1-x)/0.3)};
\addlegendentry{$n=10$, $7$ heads\quad}
\addplot[distblue, very thick] {70*ln(x/0.7)+30*ln((1-x)/0.3)};
\addlegendentry{$n=100$, $70$ heads}
\draw[gray, dashed] (axis cs:0.7,-6) -- (axis cs:0.7,0.3);
\node[gray, anchor=south, font=\small] at (axis cs:0.7,0.3) {$\hat p=0.7$};
\end{axis}
\end{tikzpicture}
\caption{Same MLE, different evidence. With ten times the data the log-likelihood is ten
times as curved, so the plausible range of $p$ is about $\sqrt{10}\approx3.2$ times
narrower.}
\label{fig:lik-curvature}
\end{figure}

\begin{definition}
The \textbf{observed information} is $J(\hat\theta) = -\ell''(\hat\theta)$, the
curvature of the log-likelihood at its peak. Its expected value over datasets,
$I(\theta) = -\E[\ell''(\theta)]$, is the \textbf{Fisher information}. Both measure how
sharply the data pin down $\theta$.
\end{definition}

\begin{theorem}[Large-sample behaviour of the MLE]
Under regularity conditions (smooth model, $\theta$ not at the edge of its range), for
large $n$
\[
  \hat\theta_{\text{MLE}} \;\overset{\text{approx}}{\sim}\;
  \Normal\!\Big(\theta,\ \frac{1}{I(\theta)}\Big),
  \qquad
  \SE(\hat\theta) \approx \frac{1}{\sqrt{I(\hat\theta)}} .
\]
\end{theorem}

\begin{whybox}
Near its peak, any smooth hill looks like a parabola:
$\ell(\theta)\approx\ell(\hat\theta) - \tfrac12 J\,(\theta-\hat\theta)^2$. That is the
log of a Normal density in $\theta$ with variance $1/J$. Curvature and variance are
reciprocals. The full proof (a Taylor expansion of $\ell'$ plus the CLT) is in
Probability From the Ground Up, Ch.~30.
\end{whybox}

\begin{example}
For the coin, $\ell''(p) = -k/p^2 - (n-k)/(1-p)^2$. At $\hat p = k/n$ this simplifies to
$-n/(\hat p(1-\hat p))$. With $7$ heads in $10$: $J = 10/0.21 = 47.6$ and
\[
  \SE(\hat p) = \sqrt{\frac{\hat p(1-\hat p)}{n}} = \sqrt{0.021} = 0.145 .
\]
This is the familiar standard error of a proportion from Chapter 7, now derived from
curvature. With $70$ heads in $100$ it is $0.046$. For a Normal mean, $I = n/\sigma^2$
gives $\SE = \sigma/\sqrt n$. For the Poisson rate, $I = n/\lambda$ gives
$\SE(\hat\lambda) = \sqrt{2/6} = 0.58$ errors per hour.
\end{example}

\begin{happeningbox}
This is how software reports standard errors for logistic regression and most fitted
models. It finds the peak of $\ell$, computes the matrix of second derivatives there
(the Hessian), inverts it, and takes square roots of the diagonal. Every ``std err''
column in a regression table is a curvature reading.
\end{happeningbox}

\subsection{Likelihood ratios}

To compare two parameter values, compare how well they predict the data. The
\textbf{likelihood-ratio statistic} for testing $H_0\colon \theta=\theta_0$ is
\[
  \Lambda = 2\big[\ell(\hat\theta) - \ell(\theta_0)\big],
\]
twice the drop in log-likelihood from the peak down to $\theta_0$. When $H_0$ is true
and $n$ is large, $\Lambda$ approximately follows a $\chi^2_1$ distribution (Chapter 10),
whose $95$th percentile is $3.84$.

For the coin and $H_0\colon p=0.5$: $\Lambda = 2(-6.109 + 6.931) = 1.65$, below
$3.84$, with $p$-value $0.20$. Seven heads in ten is not strong evidence against a fair
coin. Inverting the test gives a \textbf{likelihood interval}: all $p$ whose
log-likelihood is within $3.84/2 = 1.92$ of the peak.

\begin{figure}[H]
\centering
\begin{tikzpicture}
\begin{axis}[
  width=11cm, height=6cm, domain=0.2:0.97, samples=200,
  xlabel={$p$}, ylabel={$\ell(p) - \ell(\hat p)$},
  ymin=-5, ymax=0.7, axis lines=left,
]
\addplot[distblue, very thick] {7*ln(x/0.7)+3*ln((1-x)/0.3)};
\draw[densityred, dashed] (axis cs:0.2,-1.92) -- (axis cs:0.97,-1.92);
\node[densityred, anchor=north, font=\small] at (axis cs:0.655,-2.05) {cut at $-1.92$};
\draw[targetgreen, line width=3pt] (axis cs:0.394,-4.6) -- (axis cs:0.915,-4.6);
\node[targetgreen, anchor=south, font=\small] at (axis cs:0.655,-4.5) {95\% likelihood interval $(0.39,\ 0.92)$};
\draw[targetgreen, dotted] (axis cs:0.394,-1.92) -- (axis cs:0.394,-4.6);
\draw[targetgreen, dotted] (axis cs:0.915,-1.92) -- (axis cs:0.915,-4.6);
\fill[gray] (axis cs:0.5,-0.8228) circle (2pt);
\node[gray, anchor=north west, font=\small] at (axis cs:0.51,-0.95) {$p=0.5$: drop $0.82$};
\end{axis}
\end{tikzpicture}
\caption{The likelihood-ratio test of $p=0.5$ asks how far $\ell$ falls from its peak.
It falls only $0.82$ (so $\Lambda=1.65$), well inside the cut. The set of $p$ above the
cut is a 95\% interval; unlike $\hat p\pm1.96\,\SE$ it is lopsided, following the shape
of the hill.}
\label{fig:lik-lr}
\end{figure}

The Wald interval $0.7 \pm 1.96(0.145) = (0.42,\ 0.98)$ forces symmetry. With small $n$
and $p$ near an edge, the likelihood interval is the more trustworthy of the two. The
``deviance'' printed by logistic-regression software is the same kind of statistic: twice a
log-likelihood gap, here between a model and a smaller model nested inside it.

\begin{historybox}
In 1922 the British statistician and geneticist Ronald Fisher published ``On the
mathematical foundations of theoretical statistics''. He was stuck on a basic question:
given data and a model with unknown constants, which recipe gives the best estimate? The
competing methods of his day, such as the method of moments, gave different answers with
no principled way to choose. Fisher proposed choosing the value that maximises what he
named the \emph{likelihood}, and in the same paper defined consistency (closing in on
the truth with more data) and efficiency (having the smallest possible variance). His
later notion of \emph{information}, the curvature above, told scientists how many
observations an experiment needed. Nearly every model fitted today, from logistic
regression to a neural network, is trained by Fisher's recipe.
\end{historybox}

\section{The Bayesian view}
\label{sec:lik-bayes}

\subsection{From likelihood to a distribution over $\theta$}

Likelihood ranks parameter values but cannot say ``there is a 90\% chance $p$ exceeds
$0.5$''. Bayes' rule can, if we are willing to state what we believed before the data.

\begin{definition}
The \textbf{prior} $\pi(\theta)$ is a probability distribution over the parameter that
encodes belief before seeing the data. The \textbf{posterior} $\pi(\theta\given\text{data})$
is the belief afterwards. Bayes' rule links them:
\[
  \pi(\theta\given\text{data}) =
  \frac{L(\theta)\,\pi(\theta)}{\int L(t)\,\pi(t)\,\dd t}
  \;\propto\; \text{likelihood}\times\text{prior}.
\]
The denominator, the \textbf{evidence}, is only a normalising constant that makes the
posterior integrate to $1$.
\end{definition}

\begin{figure}[H]
\centering
\begin{tikzpicture}[
  box/.style={draw, rounded corners, minimum width=2.4cm, minimum height=0.9cm, align=center, font=\small},
  >=Stealth]
\node[box, fill=targetgreen!15] (prior) at (0,0.8) {prior $\pi(\theta)$\\ belief before};
\node[box, fill=histfill!30] (lik) at (0,-0.8) {likelihood $L(\theta)$\\ what data say};
\node[draw, circle, minimum size=0.8cm] (times) at (3,0) {$\times$};
\node[box] (norm) at (5.6,0) {rescale to\\ area $1$};
\node[box, fill=densityred!15] (post) at (8.8,0) {posterior\\ belief after};
\draw[->] (prior) -- (times); \draw[->] (lik) -- (times);
\draw[->] (times) -- (norm); \draw[->] (norm) -- (post);
\draw[->, dashed] (post.north) to[out=140, in=30] node[above, font=\scriptsize] {becomes the prior for the next batch} (prior.north);
\end{tikzpicture}
\caption{The Bayesian update. Multiply belief by evidence, rescale, and the result can
serve as the prior when more data arrive.}
\label{fig:lik-bayes-flow}
\end{figure}

\subsection{Beta--Binomial: the worked case}

For a proportion $p$ the convenient prior is the \textbf{Beta distribution}
$\Bet(a,b)$, with density proportional to $p^{a-1}(1-p)^{b-1}$ on $[0,1]$. Think of $a$
and $b$ as \textbf{pseudo-counts}: imaginary heads and tails you have ``seen'' before any
real data. Its mean is $a/(a+b)$, the pseudo-head fraction. $\Bet(1,1)$ is flat;
$\Bet(2,2)$ is a gentle hump at $0.5$; $\Bet(20,20)$ is a firm belief that $p$ is near
$0.5$.

Multiply the prior by the likelihood $p^k(1-p)^{n-k}$:
\[
  p^{a-1}(1-p)^{b-1}\cdot p^k(1-p)^{n-k} = p^{a+k-1}(1-p)^{b+n-k-1} .
\]
That is another Beta. The update is just counting:
\[
  \Bet(a,b) \;\xrightarrow{\;k \text{ heads},\ n-k \text{ tails}\;}\; \Bet(a+k,\ b+n-k).
\]
A prior that stays in the same family after updating is called \textbf{conjugate}.

\begin{example}
Prior $\Bet(2,2)$, data $7$ heads in $10$. The posterior is $\Bet(9,5)$, with
\begin{itemize}
  \item posterior mean $9/14 = 0.643$, pulled from $0.7$ toward the prior's $0.5$;
  \item posterior mode $(9-1)/(9+5-2) = 8/12 = 0.667$;
  \item posterior SD $0.124$, close to the frequentist SE of $0.145$;
  \item 95\% \textbf{credible interval} $(0.386,\ 0.861)$: the central 95\% of the
        posterior.
\end{itemize}
\end{example}

\begin{figure}[H]
\centering
\begin{tikzpicture}
\begin{axis}[
  width=11.5cm, height=6.2cm, domain=0.001:0.999, samples=200,
  xlabel={$p$}, ylabel={density},
  ymin=0, ymax=3.6, axis lines=left,
  legend style={at={(0.02,0.98)}, anchor=north west, font=\small},
]
\addplot[targetgreen, thick, dashed] {6*x*(1-x)};
\addlegendentry{prior $\Bet(2,2)$}
\addplot[distblue, thick, dotted] {1320*x^7*(1-x)^3};
\addlegendentry{likelihood (rescaled)}
\addplot[name path=post, densityred, very thick] {6435*x^8*(1-x)^4};
\addlegendentry{posterior $\Bet(9,5)$}
\addplot[draw=none, name path=base, domain=0.386:0.861] {0};
\addplot[densityred!20, domain=0.386:0.861] fill between[of=post and base, soft clip={domain=0.386:0.861}];
\node[densityred, font=\small, anchor=north] at (axis cs:0.62,0.7) {95\% credible};
\end{axis}
\end{tikzpicture}
\caption{The posterior sits between the prior and the likelihood, closer to the
likelihood because ten flips outweigh the prior's two pseudo-flips each way. The shaded
region holds 95\% of posterior probability.}
\label{fig:lik-beta-binomial}
\end{figure}

\subsection{How much does the prior matter?}

The posterior mean $(a+k)/(a+b+n)$ adds the pseudo-counts to the real counts as if
they were real flips. So a prior is worth $a+b$ observations: $\Bet(2,2)$ is worth $4$
flips, $\Bet(20,20)$ is worth $40$.

\begin{figure}[H]
\centering
\begin{tikzpicture}
\begin{axis}[
  name=weak, width=6.4cm, height=5.2cm, domain=0.001:0.999, samples=150,
  title={weak prior $\Bet(2,2)$}, title style={font=\small},
  xlabel={$p$}, ymin=0, ymax=6.2, axis lines=left,
]
\addplot[targetgreen, thick, dashed] {6*x*(1-x)};
\addplot[densityred, very thick] {6435*x^8*(1-x)^4};
\draw[gray, dotted] (axis cs:0.7,0) -- (axis cs:0.7,5.5);
\node[gray, font=\scriptsize, anchor=south] at (axis cs:0.7,5.5) {MLE};
\end{axis}
\begin{axis}[
  at={(weak.east)}, anchor=west, xshift=1.2cm,
  width=6.4cm, height=5.2cm, domain=0.01:0.99, samples=150,
  title={strong prior $\Bet(20,20)$}, title style={font=\small},
  xlabel={$p$}, ymin=0, ymax=6.2, axis lines=left,
]
\addplot[targetgreen, thick, dashed] {exp(27.952+19*ln(x)+19*ln(1-x))};
\addplot[densityred, very thick] {exp(34.8329+26*ln(x)+22*ln(1-x))};
\draw[gray, dotted] (axis cs:0.7,0) -- (axis cs:0.7,5.5);
\node[gray, font=\scriptsize, anchor=south] at (axis cs:0.7,5.5) {MLE};
\end{axis}
\end{tikzpicture}
\caption{Same data ($7$ of $10$), two priors (dashed). The weak prior lets the data
speak: posterior mean $0.64$. The strong prior, worth $40$ flips, overrules ten flips:
posterior $\Bet(27,23)$ with mean $0.54$.}
\label{fig:lik-prior-strength}
\end{figure}

As data accumulate, any reasonable prior gets swamped. Keep the observed rate at $70\%$
and grow $n$:

\begin{figure}[H]
\centering
\begin{tikzpicture}
\begin{axis}[
  width=11.5cm, height=6cm, samples=200,
  xlabel={$p$}, ylabel={posterior density},
  xmin=0, xmax=1, ymin=0, ymax=13.5, axis lines=left,
  legend style={at={(0.02,0.98)}, anchor=north west, font=\small},
]
\addplot[targetgreen, thick, dashed, domain=0.001:0.999] {6*x*(1-x)};
\addlegendentry{prior $\Bet(2,2)$}
\addplot[densityred, thick, domain=0.001:0.999] {6435*x^8*(1-x)^4};
\addlegendentry{$n=10$: $\Bet(9,5)$}
\addplot[quotientorange, thick, domain=0.3:0.98] {exp(33.9396+36*ln(x)+16*ln(1-x))};
\addlegendentry{$n=50$: $\Bet(37,17)$}
\addplot[distblue, thick, domain=0.5:0.9] {exp(126.2487+141*ln(x)+61*ln(1-x))};
\addlegendentry{$n=200$: $\Bet(142,62)$}
\end{axis}
\end{tikzpicture}
\caption{Sequential updating with 70\% heads throughout. The posterior narrows (SD
$0.124$, $0.063$, $0.032$) and its mean approaches $0.7$ ($0.643$, $0.685$, $0.696$).
With enough data, Bayesian and maximum-likelihood answers coincide.}
\label{fig:lik-sequential}
\end{figure}

\begin{recoverybox}
Pause and think: you start from $\Bet(9,5)$ and flip ten more times, getting $6$ heads.
What is the new posterior? ($\Bet(15,9)$. Updating in two batches gives the same answer
as updating once with $13$ heads in $20$ from $\Bet(2,2)$.)
\end{recoverybox}

\subsection{Credible intervals versus confidence intervals}

The two kinds of interval answer different questions.

\begin{itemize}
  \item A 95\% \textbf{confidence interval} (Chapter 7) is produced by a recipe that, over
  many repeated studies, captures the fixed true $p$ 95\% of the time. The randomness is
  in the interval. For one realised interval, ``95\% probability $p$ is inside'' is not a
  statement the method makes.
  \item A 95\% \textbf{credible interval} contains 95\% of the posterior probability.
  Given the prior and the model, ``$p$ is in here with probability 0.95'' is exactly
  what it means. The randomness is in your knowledge of $p$.
\end{itemize}

\begin{figure}[H]
\centering
\begin{tikzpicture}[x=11cm]
\draw[->] (0,0) -- (1.05,0) node[right, font=\small] {$p$};
\foreach \t in {0,0.2,0.4,0.6,0.8,1} \draw (\t,0.08) -- (\t,-0.08) node[below, font=\scriptsize] {\t};
\draw[gray, dashed] (0.7,-0.1) -- (0.7,3.3);
\node[gray, font=\scriptsize, anchor=south] at (0.7,3.3) {$\hat p = 0.7$};
\draw[sourceblue, line width=2.2pt] (0.416,2.6) -- (0.984,2.6);
\node[sourceblue, font=\small, anchor=east] at (0.40,2.6) {Wald CI};
\draw[targetgreen, line width=2.2pt] (0.394,2.0) -- (0.915,2.0);
\node[targetgreen, font=\small, anchor=east] at (0.38,2.0) {likelihood interval};
\draw[structurepurple, line width=2.2pt] (0.390,1.4) -- (0.891,1.4);
\node[structurepurple, font=\small, anchor=east] at (0.375,1.4) {credible, flat prior};
\draw[densityred, line width=2.2pt] (0.386,0.8) -- (0.861,0.8);
\node[densityred, font=\small, anchor=east] at (0.37,0.8) {credible, $\Bet(2,2)$};
\end{tikzpicture}
\caption{Four 95\% intervals for $p$ after $7$ of $10$ heads. With this little data
the method matters: the symmetric Wald interval stretches to $0.98$, while the
likelihood and Bayesian intervals respect that $p$ cannot exceed $1$ and lean left.}
\label{fig:lik-intervals}
\end{figure}

With a flat prior $\Bet(1,1)$ the posterior is $\Bet(8,4)$, interval $(0.390,\ 0.891)$.
For large $n$ and a weak prior, credible and confidence intervals nearly coincide
numerically; the difference is mostly in how you are allowed to read them.

\begin{lstlisting}[language=Python]
from scipy import stats
a, b, k, n = 2, 2, 7, 10
post = stats.beta(a + k, b + n - k)          # Beta(9, 5)
print(post.mean(), post.interval(0.95))       # 0.643, (0.386, 0.861)
print(1 - post.cdf(0.5))                      # P(p > 0.5 | data) = 0.867
\end{lstlisting}

The last line is something only the Bayesian view can give: a direct probability
that the coin favours heads.

\section{MAP estimation is regularisation}
\label{sec:lik-map}

\subsection{The posterior mode as a penalised MLE}

Often we want one number rather than a whole posterior. The \textbf{maximum a
posteriori} (MAP) estimate is the posterior's peak. Take logs of Bayes' rule; the
evidence is a constant and drops out:
\[
  \hat\theta_{\text{MAP}} = \argmax_\theta\big[\,\ell(\theta) + \log\pi(\theta)\,\big]
  = \argmin_\theta\big[\,\underbrace{-\ell(\theta)}_{\text{loss}}
  \;+\; \underbrace{(-\log\pi(\theta))}_{\text{penalty}}\,\big].
\]
MAP is the MLE plus a penalty, and the penalty is the negative log prior. A flat prior
adds nothing, so MAP equals MLE. For the coin with $\Bet(2,2)$, the MAP is $0.667$, the
posterior mode found above.

\subsection{Gaussian prior gives ridge; Laplace prior gives lasso}

Put a Normal prior $\Normal(0,\tau^2)$ on each regression weight $\beta_j$. Then
$-\log\pi(\beta) = \sum_j\beta_j^2/(2\tau^2) + \text{const}$. With Gaussian noise of
variance $\sigma^2$, multiply the whole objective by $2\sigma^2$:
\[
  \hat\beta_{\text{MAP}} = \argmin_\beta \Big[\sum_i (y_i - x_i^\top\beta)^2
  + \lambda\sum_j \beta_j^2\Big], \qquad \lambda = \frac{\sigma^2}{\tau^2}.
\]
That is \textbf{ridge regression} (L2 regularisation, or ``weight decay'' in neural
networks). A tight prior (small $\tau$) means a big $\lambda$ and strong shrinkage
toward zero.

A Laplace prior with scale $b$ gives $-\log\pi(\beta) = \sum_j\abs{\beta_j}/b$ and so
\[
  \hat\beta_{\text{MAP}} = \argmin_\beta \Big[\sum_i (y_i - x_i^\top\beta)^2
  + \lambda\sum_j \abs{\beta_j}\Big], \qquad \lambda = \frac{2\sigma^2}{b} .
\]
That is the \textbf{lasso} (L1 regularisation).

\begin{figure}[H]
\centering
\begin{tikzpicture}
\begin{axis}[
  name=prior, width=6.4cm, height=5cm, domain=-3:3, samples=161,
  title={prior on a weight}, title style={font=\small},
  xlabel={$\beta$}, ymin=0, ymax=0.95, axis lines=left,
  legend style={at={(0.02,0.98)}, anchor=north west, font=\scriptsize},
]
\addplot[distblue, thick] {exp(-x^2/2)/sqrt(2*pi)};
\addlegendentry{Gaussian}
\addplot[densityred, thick] {exp(-abs(x)/0.6)/1.2};
\addlegendentry{Laplace}
\end{axis}
\begin{axis}[
  at={(prior.east)}, anchor=west, xshift=1.4cm,
  width=6.4cm, height=5cm, domain=-3:3, samples=161,
  title={penalty $=-\log$ prior}, title style={font=\small},
  xlabel={$\beta$}, ymin=0, ymax=5, axis lines=left,
  legend style={at={(0.5,0.98)}, anchor=north, font=\scriptsize},
]
\addplot[distblue, thick] {x^2/2};
\addlegendentry{L2 (ridge)}
\addplot[densityred, thick] {abs(x)/0.6};
\addlegendentry{L1 (lasso)}
\end{axis}
\end{tikzpicture}
\caption{The Laplace prior has a sharp spike at zero; its penalty is a V with a corner
there. The Gaussian prior is smooth at zero; its penalty is a bowl. The corner is what
lets lasso set weights exactly to zero.}
\label{fig:lik-priors-penalties}
\end{figure}

\begin{example}
Estimate a single mean with $n=4$ observations, $\xbar = 6$, noise $\sigma = 2$, and
prior $\mu\sim\Normal(0,1)$. Then $\lambda = \sigma^2/\tau^2 = 4$ and the ridge
objective $\sum(x_i-\mu)^2 + 4\mu^2$ is minimised at
\[
  \hat\mu_{\text{MAP}} = \frac{n\xbar}{n+\lambda} = \frac{24}{8} = 3 .
\]
The data say $6$, the prior says $0$, and they carry equal weight here (data precision
$n/\sigma^2 = 1$ equals prior precision $1/\tau^2 = 1$), so the answer is halfway.
With a Laplace prior of scale $b=1$, the MAP is the ``soft-threshold'' of $\xbar$: move
it toward zero by $\sigma^2/(nb) = 1$, giving $5$, and stop at zero if you would cross
it.
\end{example}

\begin{figure}[H]
\centering
\begin{tikzpicture}
\begin{axis}[
  width=9cm, height=6.4cm, domain=-3:3, samples=121,
  xlabel={unpenalised estimate $\hat\beta_{\text{MLE}}$},
  ylabel={penalised estimate},
  xmin=-3, xmax=3, ymin=-3, ymax=3, axis lines=left,
  extra x ticks={0}, extra y ticks={0},
  extra tick style={grid=major, grid style={gray!40}},
  legend style={at={(0.02,0.98)}, anchor=north west, font=\footnotesize, cells={anchor=west}},
]
\addplot[gray, dashed] {x};
\addlegendentry{no penalty}
\addplot[distblue, very thick] {x/2};
\addlegendentry{ridge: scale by $\tfrac12$}
\addplot[densityred, very thick] {(x>1)*(x-1) + (x<-1)*(x+1)};
\addlegendentry{lasso: shift by $1$, clip at $0$}
\end{axis}
\end{tikzpicture}
\caption{How each penalty transforms an estimate, using the numbers of the example
above (one weight, or several whose predictors are uncorrelated). Ridge
shrinks everything proportionally and never reaches zero. Lasso subtracts a fixed
amount and sets small estimates exactly to zero, which performs variable selection.}
\label{fig:lik-shrinkage}
\end{figure}

\begin{intuitionbox}
Chapter 13 mentioned ridge and lasso as penalties added to the loss to keep
coefficients small. On its own that looks like a hack: add a term because it helps on
test data. Through the Bayesian lens it is honest modelling. You believe most
weights are small (Gaussian prior) or that most are exactly zero and a few matter
(Laplace prior), and the penalty is that belief written as a loss. Returning to
Figure~\ref{fig:lik-divisor-mse}: the prior adds bias toward zero and buys a larger cut
in variance.
\end{intuitionbox}

\begin{trapbox}
MAP is a point estimate, not a full Bayesian answer. It throws away the posterior's
width, and it depends on how you parametrise: the mode of a density changes if you
switch from $p$ to $\log p$, while posterior probabilities of events do not. If you need
uncertainty, report a credible interval, not just the MAP.
\end{trapbox}

\section{When frequentist versus Bayesian matters in practice}
\label{sec:lik-practice}

With plenty of data and a weak prior, both schools return the same estimate and nearly
the same interval, as Figure~\ref{fig:lik-sequential} showed. The choice matters at
the edges.

\begin{center}
\begin{tabular}{>{\raggedright\arraybackslash}p{4.2cm}>{\raggedright\arraybackslash}p{4cm}>{\raggedright\arraybackslash}p{4cm}}
\toprule
situation & frequentist habit & Bayesian habit \\
\midrule
large $n$, simple model & MLE $\pm 1.96\,\SE$ & essentially identical \\
small $n$ or zero events (0 of 20 failures) & MLE $=0$, Wald SE $=0$: nonsense & prior keeps estimate above $0$ \\
many related groups (per-store rates) & estimate each separately & shrink toward shared mean (hierarchical prior) \\
A/B test readout & $p$-value, fixed sample size & $\P(\text{B beats A}\given\text{data})$ \\
regularised ML models & penalty chosen by cross-validation & penalty is a prior; same fit \\
regulated settings (drug trials) & required error-rate guarantees & used with care, prior must be justified \\
\bottomrule
\end{tabular}
\end{center}

\begin{figure}[H]
\centering
\begin{tikzpicture}[
  q/.style={draw, rounded corners, fill=histfill!25, align=center, font=\small, minimum height=0.8cm, text width=3.4cm},
  a/.style={draw, rounded corners, align=center, font=\small, minimum height=0.8cm, text width=2.9cm},
  >=Stealth, node distance=0.6cm and 0.9cm]
\node[q] (q1) {Lots of data per parameter?};
\node[a, fill=targetgreen!15, right=of q1] (a1) {Either; MLE with SE is fine};
\node[q, below=of q1] (q2) {Real prior knowledge, or many similar groups?};
\node[a, fill=densityred!15, right=of q2] (a2) {Bayesian or regularised};
\node[q, below=of q2] (q3) {Need guaranteed long-run error rates?};
\node[a, fill=sourceblue!15, right=of q3] (a3) {Frequentist CI and test};
\node[a, below=of q3] (a4) {Report both; check they agree};
\draw[->] (q1) -- node[above, font=\scriptsize] {yes} (a1);
\draw[->] (q1) -- node[left, font=\scriptsize] {no} (q2);
\draw[->] (q2) -- node[above, font=\scriptsize] {yes} (a2);
\draw[->] (q2) -- node[left, font=\scriptsize] {no} (q3);
\draw[->] (q3) -- node[above, font=\scriptsize] {yes} (a3);
\draw[->] (q3) -- node[left, font=\scriptsize] {no} (a4);
\end{tikzpicture}
\caption{A rough guide. Most of the time the answers agree; the decision is about the
cases where data are thin.}
\label{fig:lik-decision}
\end{figure}

\begin{example}
A new checkout flow has $0$ failures in $20$ trials. The MLE failure rate is $0$, with a
Wald SE of $0$, an interval that claims certainty. A flat $\Bet(1,1)$ prior gives a
$\Bet(1,21)$ posterior with mean $1/22 = 0.045$ and 95\% upper bound
$1-0.05^{1/21} = 0.133$. The frequentist ``rule of three'' gives a similar upper bound,
$3/20 = 0.15$. Both say the same thing: zero failures in twenty does not mean the
rate is zero.
\end{example}

For depth, Probability From the Ground Up treats estimation theory in Ch.~30 and
Bayesian inference in Ch.~31. Chapter 15 uses this chapter's dictionary to read the loss
functions, regularisers and $\E_{x\sim p}[\cdot]$ notation of machine learning.

\section{Summary}

\begin{itemize}
  \item An \textbf{estimator} is a recipe from data to a number; it has a sampling
  distribution. $\MSE = \Bias^2 + \Var$, so a little bias can buy a big cut in variance.
  \item \textbf{Consistency}: the sampling distribution collapses onto $\theta$ as
  $n\to\infty$.
  \item The \textbf{likelihood} $L(\theta)=p(\text{data}\given\theta)$ scores parameters
  by how well they predict the data. Work with $\ell = \log L$.
  \item MLEs: Bernoulli $\hat p = k/n$; Normal $\hat\mu = \xbar$,
  $\hat\sigma^2 = SS/n$ (biased, unlike $s^2$); Poisson $\hat\lambda = \xbar$.
  \item Losses are negative log-likelihoods: squared error $\leftrightarrow$ Gaussian,
  absolute error $\leftrightarrow$ Laplace, cross-entropy $\leftrightarrow$ Bernoulli or
  categorical.
  \item Curvature of $\ell$ is information: $\SE(\hat\theta)\approx 1/\sqrt{I(\hat\theta)}$.
  \item Likelihood ratio $\Lambda = 2[\ell(\hat\theta)-\ell(\theta_0)]\approx\chi^2_1$;
  cut at $3.84$.
  \item Bayes: posterior $\propto$ likelihood $\times$ prior. Beta--Binomial:
  $\Bet(a,b)\to\Bet(a+k,\,b+n-k)$.
  \item Credible intervals are probability statements about $\theta$; confidence intervals
  are statements about the recipe.
  \item MAP $=$ MLE $+$ penalty $-\log\pi$. Gaussian prior $=$ ridge (L2), Laplace prior
  $=$ lasso (L1).
\end{itemize}

\section{Exercises}

\begin{enumerate}
  \item An estimator has bias $0.5$ and variance $2$. What is its MSE? A second estimator
  is unbiased with variance $2.4$. Which has lower MSE?
  \item A feature was used in $3$ of $12$ sessions. Find the MLE of the usage rate and its
  approximate standard error from the Fisher information.
  \item Daily incident counts over five days are $5, 3, 4, 6, 2$. Assuming a Poisson
  model, find $\hat\lambda$ and $\SE(\hat\lambda)$.
  \item For the data $2, 4, 6, 8$ compute the Normal MLE $\hat\sigma^2$ and the sample
  variance $s^2$. Explain the difference in one sentence.
  \item For $\Unif(0,\theta)$ data, show that $\hat\theta = 2\xbar$ is unbiased with
  variance $\theta^2/(3n)$. (Use $\Var(X) = \theta^2/12$.)
  \item With a flat $\Bet(1,1)$ prior and the data of Exercise 2, give the posterior, its
  mean and its mode. Which of these equals the MLE?
  \item A coin shows $60$ heads in $100$ flips. Compute the likelihood-ratio statistic for
  $H_0\colon p = 0.5$ and decide at the 5\% level.
  \item Ten observations have $\xbar = 2$. Find the ridge estimate of the mean,
  minimising $\sum(x_i-\mu)^2 + \lambda\mu^2$ with $\lambda = 10$. What prior variance
  $\tau^2$ does this correspond to if $\sigma^2 = 1$?
  \item You are training a model to predict the 99th percentile of request latency.
  Which loss from the dictionary should you use, and why not squared error?
  \item Two analysts use priors $\Bet(1,1)$ and $\Bet(30,30)$ for the same conversion rate.
  After $400$ conversions in $1000$ visits, compute both posterior means. Why are they
  close?
\end{enumerate}

\subsection*{Answers}

\begin{enumerate}
  \item $0.5^2 + 2 = 2.25$. The unbiased one has MSE $2.4$, so the biased estimator is
  better.
  \item $\hat p = 0.25$; $\SE = \sqrt{0.25\cdot0.75/12} = 0.125$.
  \item $\hat\lambda = 20/5 = 4$; $I = n/\lambda$ so $\SE = \sqrt{4/5} = 0.894$.
  \item $\xbar = 5$, $SS = 20$; $\hat\sigma^2 = 20/4 = 5$, $s^2 = 20/3 = 6.67$. The MLE
  measures spread around $\xbar$, which hugs the data, and does not correct for it.
  \item $\E[2\xbar] = 2\cdot\theta/2 = \theta$;
  $\Var(2\xbar) = 4\cdot\theta^2/(12n) = \theta^2/(3n)$.
  \item $\Bet(4,10)$; mean $4/14 = 0.286$; mode $3/12 = 0.25$. The mode (MAP) equals the
  MLE, because a flat prior adds no penalty.
  \item $\Lambda = 2[60\log(0.6/0.5) + 40\log(0.4/0.5)] = 4.03 > 3.84$; reject at 5\%
  ($p\approx0.045$).
  \item $\hat\mu = 10\cdot2/(10+10) = 1$. $\lambda = \sigma^2/\tau^2$ gives $\tau^2 = 0.1$.
  \item The pinball loss at $\tau=0.99$; its minimiser is the 99th percentile. Squared
  error is minimised by the mean.
  \item $401/1002 = 0.400$ and $430/1060 = 0.406$. A prior worth $60$ pseudo-visits is
  small next to $1000$ real ones.
\end{enumerate}
